\section*{第 \ref{ch:g5:human-reproduction} 章——人类的生殖}

\begin{solution}{exo:g5:human-reproduction:1}
来自母亲的\omterm{def:g4:animal-reproduction:sexual}{卵细胞}与来自父亲的\omterm{def:g4:animal-reproduction:sexual}{精子}\omterm{def:g5:first-look-at-cells:cell}{细胞}；
二者结合就是\omterm{def:g5:flower-to-fruit:steps}{受精}。
\end{solution}

\begin{solution}{exo:g5:human-reproduction:2}
在母亲的\omterm{prop:g5:human-reproduction:plan}{子宫}内，持续约九个月——即\omterm{def:g5:human-reproduction:pregnancy}{怀孕}。
\end{solution}

\begin{solution}{exo:g5:human-reproduction:3}
先是\omterm{def:g5:human-reproduction:pregnancy}{胚胎}，之后是\omterm{def:g5:human-reproduction:pregnancy}{胎儿}——大约两个月时，
身体的基本构造已经形成。
\end{solution}

\begin{solution}{exo:g5:human-reproduction:4}
通过\omterm{prop:g5:human-reproduction:cord}{脐带}：\omterm{def:g4:journey-of-food:digestion}{营养物质}和氧气从母亲的\omterm{prop:g4:journey-of-food:absorb}{血液}
进入婴儿体内，婴儿产生的废物则经\omterm{prop:g5:human-reproduction:cord}{脐带}
返回母体，由母体的身体处理。
\end{solution}

\begin{solution}{exo:g5:human-reproduction:5}
你自己的\omterm{prop:g5:human-reproduction:cord}{脐带}曾附着的部位，
在出生时被剪断之前。
\end{solution}

\begin{solution}{exo:g5:human-reproduction:6}
第一声啼哭使肺第一次充满空气：
在那一刻之前，它们从未\omterm{def:g4:breathing:movements}{呼吸}过——氧气一直
经\omterm{prop:g5:human-reproduction:cord}{脐带}输送。
\end{solution}

\begin{solution}{exo:g5:human-reproduction:7}
\omterm{def:g5:flower-to-fruit:steps}{受精}发生在母体内；\omterm{def:g4:animal-reproduction:ovivivi}{胎生}分娩
（婴儿在\omterm{prop:g5:human-reproduction:plan}{子宫}内生长，出生时已是活的）；
出生后以\omterm{def:g2:animals-grow-up:born}{乳汁}喂养——正是\omterm{prop:g3:sorting-living-things:groups}{哺乳动物}纲的特征。
\end{solution}

\begin{solution}{exo:g5:human-reproduction:8}
因为母亲\omterm{prop:g4:journey-of-food:absorb}{血液}中的任何物质都会经\omterm{prop:g5:human-reproduction:cord}{脐带}
进入婴儿体内——烟和酒中的有害物质也包括在内——
而\omterm{def:g5:human-reproduction:pregnancy}{胎儿}远比成人脆弱。
\end{solution}

\begin{solution}{exo:g5:human-reproduction:9}
同卵双胞胎：一个\omterm{def:g5:flower-to-fruit:steps}{受精}\omterm{def:g4:animal-reproduction:sexual}{卵细胞}分裂成两个
\omterm{def:g5:human-reproduction:pregnancy}{胚胎}。异卵（双卵）双胞胎：两个\omterm{def:g4:animal-reproduction:sexual}{卵细胞}分别
与两个\omterm{def:g4:animal-reproduction:sexual}{精子}结合，只共享\omterm{prop:g5:human-reproduction:plan}{子宫}。同卵双胞胎
外貌相似，如同复制品，因为身体由最初的
那一个\omterm{def:g5:first-look-at-cells:cell}{细胞}构成——他们起源于同一个\omterm{def:g5:first-look-at-cells:cell}{细胞}。
\end{solution}

\begin{solution}{exo:g5:human-reproduction:10}
大象处于一端，甚至更远：通常只生一个孩子，
幼体长时间无法自理，照料不仅包括喂养和保护，
还包括教导——语言、技能、上学。鳕鱼动辄产下数百万个卵，
却无人照料，与之恰成另一极端。
\end{solution}

\begin{solution}{exo:g5:human-reproduction:11}
其任务是摄入氧气、排出二氧化碳：出生前，
母亲的\omterm{prop:g4:journey-of-food:absorb}{血液}经\omterm{prop:g5:human-reproduction:cord}{脐带}承担这一工作；
从第一次\omterm{def:g4:breathing:movements}{呼吸}起，婴儿自己的肺泡接替这一任务。
啼哭来得恰是时候，正赶在旧的供应通道
被切断的那一刻——
肺必须在\omterm{prop:g5:human-reproduction:cord}{脐带}停止供血时立即开始工作。
\end{solution}
